\documentclass[11pt,article]{article}
\usepackage[utf8]{inputenc}
\usepackage[english,russian]{babel}
\usepackage{amsmath,amssymb,mathtools,physics}
\usepackage{geometry}
\geometry{a4paper,margin=2cm}

\begin{document}

\title{\textbf{Непрерывный предел многомасштабного автомата $7n$: Эмерджентность оператора Лапласа--Бельтрами}}
\author{Исследовательская группа $7n$}
\date{\today}
\maketitle

\begin{abstract}
В данном препринте исследуется математически строгий переход от дискретных операторов инцидентности циклического автомата над полем $\mathbb{Z}_{157}$ к непрерывным дифференциальным операторам математической физики. Доказано, что спектральная плотность зон деранжирования в проконечном пределе Гротендика изоморфна спектру одномерного оператора Лапласа на компактном многообразии.
\end{abstract}

\section{Дискретная модель и оператор Кирхгофа}
Рассмотрим циклический граф $\mathcal{G} = (\mathcal{V}, \mathcal{E})$ с размерностью вершины $|\mathcal{V}| = N = 157$. Пространство состояний калибруется разбиением на непересекающиеся зоны деранжирования:
\begin{equation}
\mathcal{V} = \mathcal{D}_{60} \cup \mathcal{D}_{84} \cup \mathcal{D}_{13}, \quad \mathcal{D}_i \cap \mathcal{D}_j = \emptyset
\end{equation}
Дискретный оператор Лапласа $\Delta_{\text{disc}} \in \text{Mat}_{N \times N}(\mathbb{R})$ определяется через калибровочные сдвиги инцидентности:
\begin{equation}
(\Delta_{\text{disc}} X)_i = \sum_{a \in \mathcal{D}} (X_i - X_{i-a}) = |\mathcal{D}| \cdot X_i - \sum_{a \in \mathcal{D}} X_{i-a}
\end{equation}

\section{Спектральное представление и непрерывный предел}
Применяя дискретное преобразование Фурье (DFT), перейдем к гармоникам спектральной энтропии $\mathcal{H}_{\text{spectral}}$. Собственные значения оператора $\Delta_{\text{disc}}$ принимают вид:
\begin{equation}
\lambda_\omega = 2 \sum_{a \in \mathcal{D}} \sin^2\left(\frac{\pi \omega a}{N}\right), \quad \omega \in [0, N-1]
\end{equation}
Введем масштабный шаг пространственной сетки $h$, фиксируя макроскопическую длину контура $L = N \cdot h$. В пределе $h \to 0$ и $N \to \infty$ для низших гармоник ($\omega \ll N$) разложение в ряд Тейлора дает:
\begin{equation}
\sin^2\left(\frac{\pi \omega a}{N}\right) = \left(\frac{\pi \omega a}{N}\right)^2 - \frac{1}{3}\left(\frac{\pi \omega a}{N}\right)^4 + \mathcal{O}\left(N^{-6}\right)
\end{equation}
Следовательно, непрерывный предел оператора принимает форму:
\begin{equation}
\lambda_{\text{cont}} = \lim_{h \to 0} \frac{\lambda_\omega}{h^2} = \frac{2\pi^2 \omega^2}{L^2} \sum_{a \in \mathcal{D}} a^2 = \frac{2\pi^2 \omega^2}{L^2} I_{\mathcal{D}}
\end{equation}
где $I_{\mathcal{D}} = \sum_{a \in \mathcal{D}} a^2$ --- инвариантный структурный момент инерции калибровочной зоны. Настоящее соотношение строго изоморфно спектру непрерывного дифференциального оператора Лапласа $\Delta = -\frac{d^2}{dx^2}$ с масштабным фактором $2I_{\mathcal{D}}$.

\end{document}
